% =====================================================================
% CAP. 11 - NIVEL 1
% =====================================================================
\blocofixacao

\begin{questao}[1]{}
Um capacitor ideal possui $v_C(0^-)=\SI{7}{\volt}$. No instante $t=0$ ocorre uma
comutação no circuito. Qual é o valor de $v_C(0^+)$?
\resposta{$v_C(0^+)=\SI{7}{\volt}$}
\resolucao{A tensão em um capacitor ideal não pode variar instantaneamente, pois
$i_C=C\,\mathrm{d}v_C/\mathrm{d}t$ exigiria corrente impulsiva infinita. Logo,
$v_C(0^+)=v_C(0^-)=\SI{7}{\volt}$.}
\end{questao}

\begin{questao}[1]{}
Um indutor ideal conduz $i_L(0^-)=\SI{350}{\milli\ampere}$. Uma chave muda de
posição em $t=0$. Determine $i_L(0^+)$.
\resposta{$i_L(0^+)=\SI{350}{\milli\ampere}$}
\resolucao{A corrente de um indutor ideal é contínua. Como
$v_L=L\,\mathrm{d}i_L/\mathrm{d}t$, uma mudança instantânea de corrente exigiria
tensão impulsiva infinita. Portanto $i_L(0^+)=i_L(0^-)$.}
\end{questao}

\begin{questao}[1]{}
Um circuito RC tem $R=\SI{2}{\kilo\ohm}$ e $C=\SI{50}{\micro\farad}$. Determine
a constante de tempo.
\resposta{$\tau=\SI{100}{\milli\second}$}
\resolucao{$\tau=RC=(\pot{2}{3})(\pot{50}{-6})=\SI{0.1}{\second}
=\SI{100}{\milli\second}$.}
\end{questao}

\begin{questao}[1]{}
Um circuito RL possui $L=\SI{80}{\milli\henry}$ e resistência equivalente
$R_{\mathrm{eq}}=\SI{20}{\ohm}$. Calcule $\tau$.
\resposta{$\tau=\SI{4}{\milli\second}$}
\resolucao{$\tau=L/R_{\mathrm{eq}}=0{,}08/20=\SI{0.004}{\second}
=\SI{4}{\milli\second}$.}
\end{questao}

\begin{questao}[1]{}
Um capacitor inicialmente descarregado é carregado por uma fonte CC através de
um resistor. Após um intervalo igual a uma constante de tempo, aproximadamente
qual fração da variação total de tensão já ocorreu?
\begin{alternativas}[3]
\item \SI{36.8}{\percent}
\item \SI{50}{\percent}
\item \SI{63.2}{\percent}
\item \SI{86.5}{\percent}
\item \SI{99.3}{\percent}
\end{alternativas}
\resposta{(c)}
\resolucao{Para uma carga partindo de zero,
$v_C(\tau)=V_f(1-e^{-1})\approx0{,}632V_f$. Portanto cerca de
\SI{63.2}{\percent} da variação total já ocorreu.}
\end{questao}

\begin{questao}[1]{}
Uma resposta natural de primeira ordem é proporcional a $e^{-t/\tau}$. Qual
percentual do valor inicial resta após $3\tau$?
\resposta{$e^{-3}\times100\%\approx\SI{4.98}{\percent}$}
\resolucao{$e^{-3}\approx0{,}0498$. Assim resta cerca de \SI{4.98}{\percent} do
valor inicial.}
\end{questao}

\begin{questao}[1]{}
Ao calcular a resistência equivalente vista por um capacitor para obter a
constante de tempo, como devem ser tratadas as fontes independentes ideais?
\begin{alternativas}
\item fontes de tensão abertas e fontes de corrente em curto.
\item todas as fontes abertas.
\item fontes de tensão em curto e fontes de corrente abertas.
\item todas as fontes em curto.
\item as fontes não devem ser modificadas.
\end{alternativas}
\resposta{(c)}
\resolucao{Para calcular $R_{\mathrm{eq}}$ com fontes independentes suprimidas,
uma fonte ideal de tensão é substituída por curto-circuito e uma fonte ideal de
corrente por circuito aberto. Fontes dependentes, quando existirem, permanecem
ativas.}
\end{questao}

\begin{questao}[1]{}
Um circuito ideal possui um capacitor e um indutor independentes entre si como
únicos elementos armazenadores de energia. Qual é, em geral, a ordem máxima do
circuito?
\resposta{Segunda ordem}
\resolucao{Cada variável independente de armazenamento de energia pode introduzir
um estado. Um capacitor e um indutor independentes produzem duas variáveis de
estado, caracterizando um circuito de segunda ordem.}
\end{questao}

\begin{questao}[1]{}
Para $L=\SI{0.2}{\henry}$ e $C=\SI{50}{\micro\farad}$, determine a frequência
natural não amortecida $\omega_0$.
\resposta{$\omega_0\approx\SI{316.2}{\radian\per\second}$}
\resolucao{$LC=0{,}2\cdot\pot{50}{-6}=\pot{1}{-5}$ e
$\omega_0=1/\sqrt{LC}=1/\sqrt{\pot{1}{-5}}\approx\SI{316.2}{\radian\per\second}$.}
\end{questao}

\begin{questao}[1]{}
Em um RLC em série, $R=\SI{40}{\ohm}$, $L=\SI{0.1}{\henry}$ e
$C=\SI{100}{\micro\farad}$. Classifique o amortecimento.
\resposta{Subamortecido}
\resolucao{$\alpha=R/(2L)=40/0{,}2=\SI{200}{\per\second}$ e
$\omega_0=1/\sqrt{0{,}1\cdot\pot{100}{-6}}\approx\SI{316.2}{\radian\per\second}$.
Como $\alpha<\omega_0$, a resposta é subamortecida.}
\end{questao}

\begin{questao}[1]{}
Determine a resistência crítica de um RLC em série com $L=\SI{50}{\milli\henry}$ e
$C=\SI{200}{\micro\farad}$.
\resposta{$R_{\mathrm{crit}}\approx\SI{31.6}{\ohm}$}
\resolucao{$R_{\mathrm{crit}}=2\sqrt{L/C}=2\sqrt{0{,}05/(\pot{200}{-6})}
=2\sqrt{250}\approx\SI{31.6}{\ohm}$.}
\end{questao}

\gabarito[Nível 1]
